% =============================================================================
% Chapter 9 -- Project Management and Development Evidence (owner: Haaniah)
% Every commit hash, PR number, file count and line count in this chapter was
% read out of the repository with git / find / wc on the reporting commit.
% GenAI disclosure points at docs/genai_register.md (authoritative register).
% =============================================================================
\chapter{Project Management and Development Evidence}
\label{ch:evidence}

This chapter is the evidence trail: what has been done, by whom, with which
commits, what has not, and how the work was organised, disclosed and protected.
The worksheet is explicit that commit counts alone do not demonstrate
contribution, so each member's section below names the artefacts their commits
produced and what those artefacts do.

% ---------------------------------------------------------------------------
\section{Work Completed}
\label{sec:work-completed}

\begin{longtable}{p{0.26\textwidth} p{0.20\textwidth} p{0.30\textwidth} p{0.20\textwidth}}
\caption{Work completed during iteration 1, by owner and verification source.}\label{tab:work-completed}\\
\hline\textbf{Work item} & \textbf{State} & \textbf{Artefact} & \textbf{Evidence} \\
\hline\endfirsthead
\hline\textbf{Work item} & \textbf{State} & \textbf{Artefact} & \textbf{Evidence} \\
\hline\endhead
Panel feedback captured as a tracked plan & Complete &
  Eleven comments mapped to owned actions, deliverables and a frozen-scope list
  & \texttt{TODO.md}, commit \texttt{d8b1e18} \\
Repository layout and report skeleton & Complete &
  Report template with chapter files; project README & \texttt{d8b1e18},
  \texttt{a958add} \\
Requirements baseline (use cases, FR-01--FR-29, NFR-01--NFR-14, acceptance
  checks, traceability) & Complete & Chapter~4 & \texttt{33b00e7} \\
System design baseline and twelve diagrams & Complete (diagram commit stamps
  outstanding) & Chapter~5; \texttt{diagrams/} with rendered PDFs and one build
  script & \texttt{717a7af}, \texttt{ae0a03c}, \texttt{d2d7d0c}, \texttt{7abdfc7} \\
Reviewer surface (19 routes) & Complete against mock data &
  \texttt{apps/web} & \texttt{b09e05a}, \texttt{dcd03ac}, \texttt{d535476},
  \texttt{7b411b1} \\
Design source exported and tokens ported & Complete &
  \texttt{docs/stitch-reference/screens/} (20 HTML + 19 PNG); token block in
  \texttt{globals.css} & \texttt{b09e05a} \\
GenAI use register and prompt journal & In progress (maintained continuously) &
  \texttt{docs/genai\_register.md} & \texttt{3c84428}, \texttt{a5cdf50} \\
Dataset pilot package (15 cases, 3 repositories) & Complete for the pilot size &
  \texttt{docs/dataset/} & \texttt{66e85c7} \\
Blind annotation, second annotator & Complete (1 of 2 sheets) &
  \texttt{annotator\_B\_haaniah.csv} & \texttt{8657955} \\
Integrity checking of the dataset & Complete for the pilot &
  138 checks, 0 failures & \texttt{08\_integrity\_report.md} \\
Inter-annotator agreement & Not started (blocked on annotator A) &
  \texttt{tools/agreement.py} ready & See Ch.~8, D-03 \\
Ingestion, evidence graph, rules, verifier, model server & Not started & -- &
  Ch.~6, Ch.~10 \\
CI, branch protection, container composition & Not started & -- & Ch.~8, D-02 \\
\hline\end{longtable}

% ---------------------------------------------------------------------------
\section{Work Remaining}
\label{sec:work-remaining}

Everything in the right-hand column of Chapter~8's status table is remaining
work: adapters and the reliability layer (Haaniah), the analyser, graph and rule
engine (Mishal), the packet builder, verifier, gate, model serving and the
front-end API binding (Abdullah), plus the shared items --- annotation
completion, adjudication, CI, the traceability matrix across
requirement--design--module--test, and the demonstration script. Chapter~10
sequences it; this chapter only records that none of it is being presented as
done.

% ---------------------------------------------------------------------------
\section{Individual Responsibilities}
\label{sec:individual-responsibilities}

Responsibilities are those agreed in the proposal and restated in
Chapter~3, divided by technical stream rather than by architectural layer:

\begin{longtable}{p{0.20\textwidth} p{0.44\textwidth} p{0.32\textwidth}}
\caption{Individual responsibilities and the artefacts each team member owns.}\label{tab:responsibilities}\\
\hline\textbf{Member} & \textbf{Owns} & \textbf{Independent evaluation owned} \\
\hline\endfirsthead
\hline\textbf{Member} & \textbf{Owns} & \textbf{Independent evaluation owned} \\
\hline\endhead
Haaniah Ismail & Source adapters and read-only authenticated access; the
  normalised event envelope; idempotency, checkpointing, reconciliation and
  stale-source handling; the immutable artifact and version store.
  & Duplicate and out-of-order convergence; freshness and recovery after a
    simulated outage; ingestion latency; permission and failure behaviour \\
Mishal Fatima & Code and test analyser for the declared languages; evidence
  graph and provenance; deterministic and manual linking; drift and policy
  rules; criterion invalidation and incremental dispatch.
  & Link precision and recall; graph correctness; incremental-versus-full
    equivalence; rule precision and recall \\
Abdullah Amin & Context-packet builder; local model serving and selection;
  bounded verifier loop; adaptation experiment; verdict and provenance gate;
  reviewer surface and notification routing.
  & Retrieval, verdict and grounding metrics; base-versus-adapted comparison;
    local latency and memory; task-based user evaluation \\
\hline\end{longtable}

Four responsibilities are shared by all three: consistent terminology across
proposal, report and code; construction and adjudication of the evaluation set
(the constructor does not label their own case, and nobody validates their own
subsystem's data); end-to-end integration testing; and ethical, privacy and
tool-use disclosure.

% ---------------------------------------------------------------------------
\section{Contribution of Each Team Member}
\label{sec:contributions}

The repository holds 65 commits in total, of which 28 belong to the current
tree (the project history before \texttt{066ef56} was a deliberate reset and is
summarised below). Excluding merge commits, the 28-commit window divides
$11 / 5 / 5$ between Abdullah, Mishal and Haaniah. Counts are given last,
because the artefacts are the argument.

\subsection{Abdullah Amin}

\begin{itemize}
  \item \texttt{b09e05a} -- scaffolded \texttt{apps/web}: Next.js App Router
        structure, Stitch-derived \texttt{@theme} tokens, the app shell, the
        shared UI kit, \texttt{src/lib/types.ts} (the four-verdict vocabulary)
        and the first mock modules.
  \item \texttt{dcd03ac} -- the Conflict Inbox screen (stat cards, multi-tier
        toolbar, dense findings table).
  \item \texttt{d535476} -- normalised the graph view's node/edge coordinates
        into one 0--100 space and documented the front end.
  \item \texttt{7b411b1} -- implemented all 19 routes; merged to \texttt{main}
        as pull request \#2 (\texttt{db3f301}).
  \item \texttt{d2d7d0c} -- five diagrams (layered architecture, components,
        data-flow, verifier loop, autonomy ladder) wired into Chapter~5.
  \item \texttt{7abdfc7} -- consolidated all twelve diagrams into one folder
        with one build script, and reconstructed the missing style file for
        Haaniah's three renders.
  \item \texttt{d8b1e18} -- the TODO plan and the mid-report LaTeX skeleton;
        \texttt{152b771} and \texttt{ed0c0f6} record the cross-review of the
        other member's diagrams with follow-ups listed rather than silently
        fixed.
\end{itemize}

\subsection{Mishal Fatima}

\begin{itemize}
  \item \texttt{717a7af} -- four diagrams (use case, domain model, ERD,
        requirement-to-code anchor) plus the shared TikZ style, and their
        integration into Chapters~4 and~5; merged through the same branch.
  \item \texttt{33b00e7} -- Chapter~4: roles, six detailed use cases,
        FR-01--FR-29, NFR-01--NFR-14, per-requirement acceptance checks and the
        traceability table.
  \item \texttt{66e85c7} -- the complete dataset pilot package: repository
        screening, case schema, annotation manual, 15 case packets with head
        patches and test logs, two annotation sheets, adjudication template,
        split manifest, integrity script and dataset card.
  \item \texttt{3c84428} -- created the GenAI use register (pull request \#1,
        \texttt{74666ea}) and \texttt{af1f532} updated the team plan.
  \item Merged pull requests \#1 (\texttt{74666ea}) and \#5 (\texttt{d7cd0e2}).
\end{itemize}

\subsection{Haaniah Ismail}

\begin{itemize}
  \item \texttt{ae0a03c} -- three diagrams (ingestion sequence, event/finding
        state machines, deployment footprint).
  \item \texttt{8657955} -- completed the second blind annotation sheet: 15
        cases with verdict, freshness, review flag, reason code, evidence spans,
        justification and time per case.
  \item \texttt{0def7bb} -- updated the proposal document; \texttt{163222a} and
        \texttt{9aacf5f} record the early OAuth and Figma integration attempts.
  \item \texttt{a5cdf50} -- maintained the GenAI register; \texttt{7d3c99c}
        brought Chapters~1, 5 and the Chapter~3 draft into the repository.
  \item Merged pull requests \#3 (\texttt{afe3fd9}), \#6 (\texttt{1acd202}) and
        \#7 (\texttt{671f3a2}).
\end{itemize}

Before the reset, the shared history holds the proposal and literature-review
work of all three members, from \texttt{ae24733} (26 July 2026) onwards.

% ---------------------------------------------------------------------------
\section{Repository and Version Control Evidence}
\label{sec:repo-evidence}

\begin{longtable}{p{0.30\textwidth} p{0.62\textwidth}}
\caption{Repository and version control evidence collected from the local Git history.}\label{tab:repo-evidence}\\
\hline\textbf{Item} & \textbf{Value} \\
\hline\endfirsthead
\hline\textbf{Item} & \textbf{Value} \\
\hline\endhead
Total commits & 65 (\texttt{git rev-list --count HEAD}) \\
Commits in the current tree & 28 since \texttt{066ef56} (6 October 2026,
  ``start fresh'') \\
History span & 26 July 2026 (\texttt{ae24733}) to 9 October 2026
  (\texttt{671f3a2}, HEAD) \\
Merged pull requests & \#1, \#2, \#3, \#5, \#6, \#7, recorded in merge-commit
  messages \textbf{[TODO: verify --- pull request \#4 does not appear in local
  history; confirm on the remote before quoting a total]} \\
Remote branches & \texttt{main}, \texttt{add-genai-register},
  \texttt{feat/frontend-scaffold}, \texttt{label\_H} \\
Review practice & Feature branches merged through pull requests; two members
  cross-review the other's chapter/diagram work (recorded in
  \texttt{TODO.md}) \textbf{[TODO: verify --- branch protection and the
  who-reviews-whom rule are not configured yet]} \\
Large or late commits & The reset commit \texttt{066ef56} deleted the previous
  prototype tree; \texttt{7d3c99c} adds 1{,}004 lines of report chapters in one
  commit. Both are declared here rather than left to be discovered \\
Integration responsibility & Abdullah integrated the front end and the diagram
  set; Mishal merged the dataset and register work; Haaniah merged the label and
  report branches \\
\hline\end{longtable}

% ---------------------------------------------------------------------------
\section{Milestones and Commit Evidence}
\label{sec:milestones}

\begin{longtable}{p{0.30\textwidth} p{0.14\textwidth} p{0.24\textwidth} p{0.24\textwidth}}
\caption{Milestone and commit evidence for iteration 1.}\label{tab:milestones}\\
\hline\textbf{Milestone} & \textbf{State} & \textbf{Commit / PR} & \textbf{Artefact} \\
\hline\endfirsthead
\hline\textbf{Milestone} & \textbf{State} & \textbf{Commit / PR} & \textbf{Artefact} \\
\hline\endhead
Plan and report skeleton & Complete & \texttt{d8b1e18} & \texttt{TODO.md},
  \texttt{main.tex} \\
Proposal and document set in the repository & Complete & \texttt{a958add} &
  \texttt{docs/*.docx} \\
Diagram set (twelve) & Complete, stamps outstanding & \texttt{717a7af},
  \texttt{ae0a03c}, \texttt{d2d7d0c}, \texttt{7abdfc7} & \texttt{diagrams/} \\
Requirements chapter & Complete & \texttt{33b00e7} & Chapter~4 \\
Design chapter & Complete & \texttt{7d3c99c} (with earlier diagram commits) &
  Chapter~5 \\
Frontend scaffold + all screens & Complete (mock data) & \texttt{b09e05a} $\rightarrow$
  \texttt{7b411b1}, PR \#2 & \texttt{apps/web} \\
GenAI register & Maintained & \texttt{3c84428}, \texttt{a5cdf50} &
  \texttt{docs/genai\_register.md} \\
Dataset pilot & Complete & \texttt{66e85c7}, PR \#5 & \texttt{docs/dataset/} \\
Annotation round 1 of 2 & Partial & \texttt{8657955} & annotator B sheet \\
Local-inference spike & Not started & -- & -- \\
Ingestion / graph / verifier & Not started & -- & -- \\
\hline\end{longtable}

% ---------------------------------------------------------------------------
\section{Tools and Tool-Assisted Development}
\label{sec:tools-and-genai}

Generative AI was used and is disclosed per artefact in
\texttt{docs/genai\_register.md}, which holds the register table (date, member,
tool/model, purpose, artefact, output used, the student's own contribution, how
it was verified, any rejected or corrected output, evidence link) and a prompt
journal for the interactions that materially changed the project's direction.
The register classifies each artefact at one of the worksheet's four levels:
L1 supportive, L2 assisted, L3 core assist, L4 substituted. No artefact is
recorded at L4, and none is presented as independently completed work that an
assistant produced.

Illustrative rows already in the register:

\begin{itemize}
  \item \textbf{L3} -- Chapter~4 was drafted with an assistant from named
        proposal sections, with verification and the rewrite in the member's own
        words listed as pending student work.
  \item \textbf{L3} -- the dataset pilot (screening, case packets, manual,
        scripts) was produced with an assistant running each repository's own
        tests for every pass/fail claim; the integrity script's 138 checks and
        the correction of one wrongly claimed test result are recorded there.
  \item \textbf{L3} -- the four TikZ diagrams, with the layout corrections the
        member made listed as the student contribution.
  \item \textbf{L2} -- the blind annotation sheet, where the member made the
        case-level judgements and the assistant validated CSV structure.
  \item \textbf{L1} -- explanatory sessions used to prepare for defending work.
\end{itemize}

Non-GenAI tools in daily use: git and GitHub for version control, Google Stitch
(through its MCP integration) for screen design and export, OpenCode as the
agent harness, VS Code, Python 3.12 with pytest for the dataset scripts, Node.js
with Next.js, TypeScript, Tailwind and ESLint for the front end, and LaTeX with
the \texttt{fast-nuces-bs} class for this report. Citation checks are recorded
against the sources in \texttt{references/references.bib}.

% ---------------------------------------------------------------------------
\section{External Components and Dependencies}
\label{sec:external-components}

\begin{longtable}{p{0.24\textwidth} p{0.24\textwidth} p{0.26\textwidth} p{0.22\textwidth}}
\caption{External components and dependencies used by the project.}\label{tab:external-dependencies}\\
\hline\textbf{Dependency} & \textbf{Purpose} & \textbf{Validation performed} &
  \textbf{Risk / fallback} \\
\hline\endfirsthead
\hline\textbf{Dependency} & \textbf{Purpose} & \textbf{Validation performed} &
  \textbf{Risk / fallback} \\
\hline\endhead
Next.js 16.4.0, React, TypeScript, Tailwind, ESLint & Front end & typecheck,
  lint and production build all pass & Low; versions pinned in
  \texttt{package.json} \\
Node.js v24.14.1 & Toolchain & Build run under it & Low; advisory deprecation
  warning tracked as D-05 \\
\texttt{@material-symbols/font-400} and self-hosted fonts & Icons and type &
  Served locally; no external font requests & Low; replaceable \\
Google Stitch & Screen design and export & Exports vendored into
  \texttt{docs/stitch-reference/} so the design is reproducible without the
  service & Low; design source committed \\
GitHub, Jira, Figma APIs & Planned data sources (read-only) &
  \textbf{Not yet exercised} & Medium; documented fallback is exported/sample
  data plus a manually built Figma reference set \\
Slack & Planned notification delivery only & -- & Low; delivery is not an input
  to reasoning \\
PostgreSQL, Redis, Docker Compose & Planned deployment footprint & -- &
  Not yet adopted; single-host design \\
Open-weight model (Qwen/Llama candidates) & Planned local verifier &
  \textbf{No model downloaded or measured} & High; the hardware risk is open
  (Ch.~8, Ch.~10) \\
react-hook-form, excalidraw, mantine (all MIT) & Dataset pilot sources, pinned
  by screened commit & Install and targeted test runs recorded per repository in
  \texttt{01\_repository\_catalogue.csv} & Low; permissive licences, provenance
  recorded \\
\hline\end{longtable}

No credentials, tokens or private raw content are committed: \texttt{.env} is
gitignored, and the Stitch key is read from the environment.

% ---------------------------------------------------------------------------
\section{Risks and Mitigation}
\label{sec:risks}

The proposal's risk register is reproduced here with its current status, because
three of the entries have moved since it was written.

\begin{longtable}{p{0.34\textwidth} p{0.14\textwidth} p{0.48\textwidth}}
\caption{Risks open at the mid evaluation and the mitigation applied to each.}\label{tab:risks-current}\\
\hline\textbf{Risk} & \textbf{Status now} & \textbf{Mitigation} \\
\hline\endfirsthead
\hline\textbf{Risk} & \textbf{Status now} & \textbf{Mitigation} \\
\hline\endhead
Contemporary product judged to already cover the workflow &
  Open; narrowed & Local execution, versioned provenance, gating and abstention
    are the conceded differentiators; if a controlled comparison shows no gap,
    the claim is narrowed and reported honestly (Ch.~2) \\
Direct product trials not yet run & \textbf{Materialised as a defect} &
  Dated trial notes are outstanding; marked \textbf{[TODO: verify]} in Ch.~2 and
  scheduled before submission \\
GitHub/Jira/Figma access blocked or rate-limited & Open & Fallback to exported
  or sample data plus a manually built Figma reference set; no live access has
  been attempted yet \\
Criterion-to-code retrieval recall too low & Open (unmeasurable yet) &
  Restrict to the declared repository class and changed-symbol neighbourhood;
  report Recall@k honestly \\
Verifier produces fluent but ungrounded verdicts & Open (verifier not built) &
  Deterministic gate by construction; unsupported-assertion rate is a measured
  target of at most 2\% \\
Local inference exceeds available hardware & \textbf{Still unmeasured} &
  Smaller or quantised checkpoint; asynchronous rather than in-review analysis;
  declare the hardware rather than claim general performance \\
Gold oracle too small for fine-tuning & Open & Public component corpora for
  training, gold set reserved for evaluation; a base-model pipeline is an
  acceptable outcome \\
Team underestimates the reliability engineering & \textbf{Realised} &
  No ingestion work has started; resequenced into the next iteration with the
  scope of the first release held to Rules 1--2 \\
Scope creep toward extra integrations or visual fidelity & Controlled &
  The frozen list in \texttt{TODO.md} \S9 records what will not be built \\
\hline\end{longtable}

% ---------------------------------------------------------------------------
\section{Security and Privacy Considerations}
\label{sec:security-privacy}

\textbf{What is true today.} The deployed artefact is a static front end. It
makes no outbound requests: fonts and icons are served from the application
itself, avatars are initials, thumbnails are gradients, and no image or content
delivery network is contacted. Repository secrets are kept out of git
(\texttt{.env} is ignored), and the only secret in use, the Stitch API key, is
read from the environment. Mock data in \texttt{src/lib/mock/} contains no
credentials and no real customer content.

\textbf{What is designed but not yet implemented.} Read-only, least-privilege
OAuth scopes for the three source systems; tokens encrypted at rest and held
only by one service; artifact text treated as untrusted input, including by the
verifier, which must never act on instructions found inside issue text or commit
messages; export of audit evidence with tokens and private raw content stripped;
and egress inspection during testing to demonstrate that no repository, ticket
or design content leaves the deployment. These are NFR-09 to NFR-11 in
Chapter~4 with acceptance checks attached; none can be tested until the adapters
exist, which is stated rather than assumed.

\textbf{Position on model access.} All committed inference is local. Hosted
models appear only as an offline comparison on sanitised data and are never a
runtime dependency; this is what makes the privacy claim structural instead of
policy-based.

% ---------------------------------------------------------------------------
\section{Current Limitations}
\label{sec:limitations}

\begin{enumerate}
  \item The reviewer surface reads committed mock data; there is no API client
        and no server behind any screen.
  \item Nothing has been ingested: adapters, normalisation, idempotency,
        checkpoints, reconciliation and stale-source handling do not exist.
  \item The evidence graph, linking, rules and incremental dispatch do not
        exist, so no deterministic finding can be produced.
  \item The context packet, verifier and verdict gate do not exist; no model has
        been downloaded, and the hardware question is unanswered.
  \item There is no CI, no container composition and no automated test suite for
        the front end; correctness rests on types, lint and the build.
  \item One Python test module in the tree cannot be collected (D-01), and the
        second annotation sheet is outstanding (D-03), so agreement statistics
        are pending.
  \item The dataset is a 15-case pilot from three repositories. It is valid
        evidence that the protocol works; it is not yet an evaluation oracle.
  \item Product capabilities in Chapter~2 are read from documentation; the dated
        hands-on trials the proposal committed to have not been run.
\end{enumerate}
